\documentclass[journal]{IEEEtran}
\usepackage[utf8]{inputenc}
% babel-spanish no disponible en este entorno; tildes via inputenc utf8
% \usepackage[spanish]{babel}
\usepackage{amsmath,amssymb,amsfonts}
\usepackage{graphicx}
\usepackage{cite}
\usepackage{url}

% ============================================================
% Entrega Inicial – Teoría de Grafos 2026-I
% Secciones: Introducción, Descripción del Problema,
%            Objetivo General y Objetivos Específicos,
%            Bibliografía
% ============================================================

\begin{document}

\title{Eficiencia Espectral de la División del Trabajo Humano:\\ %
       Comparación de Particiones Emergentes con Bisecciones de Fiedler}

\author{\IEEEauthorblockN{Thomas Chísica Londoño}
\IEEEauthorblockA{Programa de Matemáticas Aplicadas y Ciencias de la Computación\\
Universidad del Rosario, Bogotá, Colombia\\
Teoría de Grafos 2026-I \quad Profesor: Daniel Alfonso Bojacá Torres\\
\texttt{thomas.chisica@urosario.edu.co}}}

\maketitle

\begin{abstract}
Este trabajo estudia si las divisiones del trabajo espacial que emergen espontáneamente
entre pares de agentes humanos que actúan sin comunicación explícita son
\emph{espectralmente eficientes}. Modelamos el espacio de interacción como un grafo
$G=(V,E)$ de $64$ nodos con 8-conectividad y utilizamos la bisección de Fiedler como
referente teórico de optimalidad. La hipótesis central es que las díadas con divisiones
focales (izquierda-derecha o arriba-abajo) convergen a particiones cuya conductancia es
comparable a la de la bisección espectral óptima $h(S_{\text{Fiedler}})$, mientras que
las díadas con divisiones mixtas exhiben ineficiencia espectral.
\end{abstract}

\begin{IEEEkeywords}
teoría algebraica de grafos, vector de Fiedler, conductancia, desigualdad de Cheeger,
coordinación multi-agente, división del trabajo, bisección espectral.
\end{IEEEkeywords}


% ============================================================
\section{Introducción}
% ============================================================

La coordinación entre agentes que no pueden comunicarse explícitamente es un
problema fundamental en sistemas multi-agente, ciencias cognitivas y teoría de
juegos~\cite{mesbahi2010}. Cuando múltiples individuos deben colaborar para cubrir
un espacio compartido sin negociación directa, emergen patrones de división del trabajo
que organizan el sistema de manera espontánea. Este fenómeno---conocido como
\emph{auto-organización}---ha sido documentado tanto en colonias de insectos como en
grupos humanos~\cite{goldstone2024}. Comprenderlo matemáticamente exige herramientas
que capturen la estructura geométrica del espacio de interacción y la calidad de sus
posibles particiones.

La teoría algebraica de grafos proporciona exactamente esas herramientas. El espectro
del operador Laplaciano de un grafo codifica propiedades globales de conectividad y,
a través de la \emph{Desigualdad de Cheeger}~\cite{chung1997,cheeger1970}, vincula el
segundo eigenvalor más pequeño $\lambda_2$ con la conductancia isoperimétrica $h(G)$.
El vector de Fiedler~\cite{fiedler1973}---el eigenvector asociado a $\lambda_2$---
produce una bisección del grafo que minimiza la conductancia entre todas las particiones
de mitad de tamaño. Este resultado es la base del \emph{spectral partitioning} y del
\emph{spectral clustering}~\cite{vonluxburg2007}, y constituye el benchmark teórico
de optimalidad que se empleará en este trabajo.

El punto de partida empírico es el experimento ``Seeking the Unicorn'' de
Andrade-Lotero y Goldstone~\cite{andrade2021}, en el cual parejas de participantes
humanos buscan un objetivo oculto en un tablero compartido de $8\times 8$ celdas sin
comunicación explícita. A lo largo de sesenta rondas, las díadas desarrollan divisiones
espaciales emergentes---típicamente izquierda-derecha o arriba-abajo---que minimizan la
redundancia y maximizan la cobertura conjunta. La pregunta que motiva este proyecto es:
\textbf{¿son esas divisiones espectralmente eficientes?}

Este trabajo no busca re-demostrar la Desigualdad de Cheeger sino \emph{usarla} como
lupa analítica: compararemos la conductancia de las particiones humanas observadas
$h(S_{\text{obs}})$ contra la conductancia de la bisección de Fiedler
$h(S_{\text{Fiedler}})$ y contra la cota inferior teórica $\lambda_2/2$. El análisis
aportará evidencia sobre si la coordinación implícita humana converge a soluciones
cercanas al óptimo espectral.


% ============================================================
\section{Descripción del Problema}
% ============================================================

\subsection{Contexto experimental}

El conjunto de datos \emph{Self-Organized Division of Cognitive Labor}
(SODCL)~\cite{andrade2021} registra el comportamiento de 45 díadas durante 60 rondas
cada una. En cada ronda, cada jugador selecciona hasta 32 celdas de un tablero de
$8\times 8$ para ``cubrir''. Las celdas visitadas simultáneamente por ambos jugadores
representan redundancia; las no visitadas por ninguno, pérdida de cobertura. Sin
instrucción explícita ni canal de comunicación, las díadas convergen a divisiones
espaciales auto-organizadas.

Un análisis exploratorio preliminar ilustra la nitidez del fenómeno: la díada
435\nobreakdash-261 exhibe una divergencia izquierda-derecha del $93.07\,\%$ en las
rondas tardías. El Jugador\,1 acumula $887$ visitas en las columnas 5--8 frente a $43$
en las columnas 1--4; el Jugador\,2 muestra el patrón exactamente invertido ($889$
versus $21$). Esta auto-segregación constituye el objeto de análisis espectral del
presente proyecto.

\subsection{Formulación graph-teórica}

\textbf{Definición 1 (Grafo del tablero).}
Sea $G=(V,E)$ el grafo donde $V=\{(i,j):1\le i,j\le 8\}$ representa las $64$ celdas
del tablero y $E$ contiene todas las aristas entre celdas adyacentes bajo
\mbox{8-conectividad} (horizontal, vertical y diagonal). Se tiene $|V|=64$ y
$|E|=224$.

\textbf{Definición 2 (Matriz Laplaciana).}
Sea $A$ la matriz de adyacencia de $G$ y $D=\operatorname{diag}(d_1,\ldots,d_n)$ la
matriz de grados. La \emph{matriz Laplaciana} es $L=D-A$. Sus eigenvalores satisfacen
$0=\lambda_1\le\lambda_2\le\cdots\le\lambda_n$.

\textbf{Definición 3 (Valor y vector de Fiedler).}
El segundo eigenvalor más pequeño de $L$, denotado $\lambda_2$, se denomina
\emph{conectividad algebraica} o \emph{valor de Fiedler}. El eigenvector asociado
$\mathbf{v}_2$ es el \emph{vector de Fiedler}.

\textbf{Definición 4 (Conductancia de una partición).}
Para una bipartición $(S,\bar{S})$ de $V$, la conductancia se define como
\begin{equation}
  h(S) \;=\;
  \frac{\bigl|\{(u,v)\in E : u\in S,\,v\in\bar{S}\}\bigr|}
       {\min\!\bigl(\operatorname{vol}(S),\,\operatorname{vol}(\bar{S})\bigr)},
\end{equation}
donde $\operatorname{vol}(S)=\sum_{v\in S}d_v$ es el volumen del conjunto $S$.

\textbf{Teorema (Desigualdad de Cheeger~\cite{chung1997}).}
Para cualquier grafo $G$,
\begin{equation}
  \frac{\lambda_2}{2} \;\le\; h(G) \;\le\; \sqrt{2\lambda_2},
\end{equation}
donde $h(G)=\min_S h(S)$ es la constante isoperimétrica del grafo.

\subsection{El problema central}

Dadas las particiones observadas $S_{\text{obs}}^d$ de cada díada $d$ y la bisección
espectral $S_{\text{Fiedler}}$ obtenida a partir del vector de Fiedler de $G$, el
problema que se aborda es el siguiente:

\begin{enumerate}
  \item \textbf{Eficiencia espectral:} ¿Qué tan cerca se encuentra $h(S_{\text{obs}}^d)$
        de la cota inferior teórica $\lambda_2/2$, que representa el límite de optimalidad
        garantizado por el Teorema de Cheeger?

  \item \textbf{Comparación con Fiedler:} ¿Es $h(S_{\text{obs}}^d)\approx
        h(S_{\text{Fiedler}})$? Es decir, ¿los humanos alcanzan particiones
        comparables a la bisección espectral computada explícitamente?

  \item \textbf{Diferencia por tipo de división:} ¿Las divisiones focales
        (izquierda-derecha o arriba-abajo) son más eficientes espectralmente que las
        divisiones mixtas?
\end{enumerate}

La relevancia del problema radica en que, si la coordinación implícita humana converge
a soluciones cercanas al óptimo espectral, ello sugeriría que los mecanismos cognitivos
de auto-organización responden, de forma no deliberada, a la estructura algebraica
subyacente del espacio de interacción.


% ============================================================
\section{Objetivo General y Objetivos Específicos}
% ============================================================

\subsection{Objetivo general}

Determinar si las divisiones del trabajo espacial que emergen espontáneamente en díadas
humanas sin comunicación son espectralmente eficientes, mediante la comparación entre la
conductancia de las particiones observadas, la bisección de Fiedler y la cota de Cheeger
del grafo del tablero experimental.

\subsection{Objetivos específicos}

\begin{enumerate}

  \item \textbf{Modelar} el tablero experimental $8\times 8$ como un grafo
        $G=(V,E)$ con 8-conectividad, construir su matriz Laplaciana $L=D-A$ y
        verificar sus propiedades algebraicas fundamentales (simetría, semidefinición
        positiva, eigenvalor nulo).

  \item \textbf{Calcular} el valor de Fiedler $\lambda_2$ y el vector de Fiedler
        $\mathbf{v}_2$ del grafo $G$, y derivar la bisección espectral
        $S_{\text{Fiedler}}=\{v: v_2(v)\ge 0\}$ junto con su conductancia
        $h(S_{\text{Fiedler}})$.

  \item \textbf{Extraer} las particiones observadas $S_{\text{obs}}^d$ para cada
        díada $d$ del conjunto de datos SODCL, a partir de las frecuencias de visita
        acumuladas por jugador en las rondas de convergencia.

  \item \textbf{Comparar} la conductancia de las particiones humanas
        $h(S_{\text{obs}}^d)$ con la conductancia espectral $h(S_{\text{Fiedler}})$ y
        con la cota inferior $\lambda_2/2$, cuantificando la eficiencia relativa
        mediante el ratio $\eta_d = h(S_{\text{Fiedler}})/h(S_{\text{obs}}^d)$.

  \item \textbf{Analizar} si existe diferencia estadísticamente significativa en la
        eficiencia espectral entre díadas con divisiones focales (izquierda-derecha,
        arriba-abajo) y díadas con divisiones mixtas.

  \item \textbf{Implementar} en Python todos los algoritmos involucrados (construcción
        del Laplaciano, descomposición espectral, bisección de Fiedler, cálculo de
        conductancia mediante recorrido en grafos) de manera reproducible y documentada.

\end{enumerate}


% ============================================================
\bibliographystyle{IEEEtran}
\begin{thebibliography}{10}

\bibitem{andrade2021}
E.\ Andrade-Lotero and R.\ L.\ Goldstone,
``Self-organized division of cognitive labor,''
\textit{PLOS ONE}, vol.~16, no.~7, p.~e0254532, Jul.\ 2021.

\bibitem{chung1997}
F.\ R.\ K.\ Chung,
\textit{Spectral Graph Theory}.
Providence, RI: American Mathematical Society, 1997.

\bibitem{fiedler1973}
M.\ Fiedler,
``Algebraic connectivity of graphs,''
\textit{Czechoslovak Mathematical Journal}, vol.~23, no.~2, pp.~298--305, 1973.

\bibitem{cheeger1970}
J.\ Cheeger,
``A lower bound for the smallest eigenvalue of the Laplacian,''
in \textit{Problems in Analysis}, R.\ C.\ Gunning, Ed.
Princeton, NJ: Princeton University Press, 1970, pp.~195--199.

\bibitem{vonluxburg2007}
U.\ von Luxburg,
``A tutorial on spectral clustering,''
\textit{Statistics and Computing}, vol.~17, no.~4, pp.~395--416, Dec.\ 2007.

\bibitem{mesbahi2010}
M.\ Mesbahi and M.\ Egerstedt,
\textit{Graph Theoretic Methods in Multiagent Networks}.
Princeton, NJ: Princeton University Press, 2010.

\bibitem{mohar1991}
B.\ Mohar,
``The Laplacian spectrum of graphs,''
in \textit{Graph Theory, Combinatorics, and Applications}, vol.~2,
Y.\ Alavi et al., Eds.\ New York: Wiley, 1991, pp.~871--898.

\bibitem{newman2018}
M.\ E.\ J.\ Newman,
\textit{Networks}, 2nd~ed.
Oxford: Oxford University Press, 2018.

\bibitem{goldstone2024}
R.\ L.\ Goldstone, E.\ Andrade-Lotero, R.\ D.\ Hawkins, and M.\ E.\ Roberts,
``The emergence of specialized roles within groups,''
\textit{Topics in Cognitive Science}, 2024.

\bibitem{spielman2007}
D.\ A.\ Spielman,
``Spectral graph theory and its applications,''
in \textit{Proc.\ 48th Annu.\ IEEE Symp.\ Foundations Comput.\ Sci.\ (FOCS)},
Providence, RI, 2007, pp.~29--38.

\end{thebibliography}

\end{document}
